\chapter{Maximum Likelihood Estimation}
\label{ch:mle}

\noindent\textsc{Almost every estimator} you'll encounter past a basic stats course is, in some sense, a maximum likelihood estimator. OLS is MLE for the normal linear model. Logistic regression is MLE for binary outcomes. Poisson regression, survival models, time-series models, almost everything: all MLE.

The principle is one of the most powerful ideas in statistics. Once you have it, derivations that previously seemed magical become almost mechanical.

\section{The likelihood function}

You have data $\bm{x} = (x_1, \ldots, x_n)$ and a model that says these are independent draws from a distribution with PDF (or PMF) $f(x; \theta)$ depending on a parameter $\theta$. The \emph{likelihood} of $\theta$ given the data is
\begin{equation}
  L(\theta \mid \bm{x}) = \prod_{i=1}^{n} f(x_i; \theta).
  \label{eq:likelihood-defn}
\end{equation}
This is the joint probability of the observed data, viewed as a function of the unknown $\theta$.

Note: the likelihood is not a probability distribution over $\theta$. It is a function of $\theta$ for fixed data. Two values of $\theta$ can be compared by their likelihood, but $L$ does not integrate to 1 over $\theta$.

\paragraph{The log-likelihood.} Products are awkward to differentiate. Logarithms turn products into sums:
\begin{equation}
  \ell(\theta) = \log L(\theta) = \sum_{i=1}^{n} \log f(x_i; \theta).
  \label{eq:log-likelihood-mle}
\end{equation}
Maximizing $L$ and maximizing $\ell$ are equivalent (log is monotonic), so we always work with $\ell$.

\section{The MLE principle}

The maximum likelihood estimator is
\begin{equation}
  \hat\theta_{\text{MLE}} = \arg\max_\theta \, \ell(\theta).
\end{equation}
Pick the parameter value that makes the observed data most plausible.

This is not a Bayesian statement. We are not saying $\theta$ is random; we are not putting a prior on it. We are saying: of all values of $\theta$, which one makes the data we actually observed most likely to occur? That value is our estimate.

The first-order condition is
\begin{equation}
  \frac{d\ell}{d\theta}(\hat\theta) = 0.
  \label{eq:score-equation}
\end{equation}
The derivative of the log-likelihood is called the \emph{score function}. The MLE is a zero of the score.

\section{Worked example: the normal mean}

$X_1, \ldots, X_n \sim N(\mu, \sigma^2)$, with both $\mu$ and $\sigma^2$ unknown. The PDF is
\[
  f(x; \mu, \sigma^2) = \frac{1}{\sqrt{2\pi\sigma^2}} \exp\!\left(-\frac{(x - \mu)^2}{2\sigma^2}\right).
\]
Log:
\[
  \log f(x; \mu, \sigma^2) = -\frac{1}{2}\log(2\pi) - \frac{1}{2}\log\sigma^2 - \frac{(x - \mu)^2}{2\sigma^2}.
\]
The full log-likelihood:
\begin{align}
  \ell(\mu, \sigma^2)
    &= -\frac{n}{2}\log(2\pi) - \frac{n}{2}\log\sigma^2 - \frac{1}{2\sigma^2}\sum_{i=1}^{n}(x_i - \mu)^2.
\end{align}

\paragraph{First-order condition for $\mu$.}
\[
  \frac{\partial \ell}{\partial \mu} = \frac{1}{\sigma^2}\sum_{i=1}^{n}(x_i - \mu) = 0 \implies \hat\mu = \bar{x}.
\]
The MLE for $\mu$ is the sample mean. This holds regardless of the value of $\sigma^2$.

\paragraph{First-order condition for $\sigma^2$.} Differentiating with respect to $\sigma^2$ (treating $\sigma^2$ as a single variable):
\[
  \frac{\partial \ell}{\partial \sigma^2} = -\frac{n}{2\sigma^2} + \frac{1}{2\sigma^4}\sum_{i=1}^{n}(x_i - \mu)^2 = 0.
\]
Solve for $\sigma^2$, plugging in $\hat\mu = \bar{x}$:
\[
  \hat\sigma^2_{\text{MLE}} = \frac{1}{n}\sum_{i=1}^{n}(x_i - \bar{x})^2.
\]
Notice the $1/n$, not $1/(n-1)$. The MLE for $\sigma^2$ is biased (it slightly underestimates the population variance). The standard ``sample variance'' you've been computing is the bias-corrected version. We'll discuss this trade-off shortly.

\section{Worked example: OLS as MLE}

Now the deeper claim: under normal errors, OLS \emph{is} the MLE for regression coefficients.

Model: $y_i = \bbeta^\top \bm{x}_i + \varepsilon_i$ with $\varepsilon_i \sim N(0, \sigma^2)$ independent. Then $y_i \mid \bm{x}_i \sim N(\bbeta^\top \bm{x}_i, \sigma^2)$. The log-likelihood:
\[
  \ell(\bbeta, \sigma^2) = -\frac{n}{2}\log(2\pi\sigma^2) - \frac{1}{2\sigma^2}\sum_{i=1}^{n}(y_i - \bbeta^\top \bm{x}_i)^2.
\]
The first term doesn't involve $\bbeta$. The second term is, up to a positive constant, the \emph{negative} of the sum of squared residuals. Maximizing $\ell$ with respect to $\bbeta$ is equivalent to minimizing $\sum (y_i - \bbeta^\top \bm{x}_i)^2$. That's OLS.

So under the normal-errors assumption, OLS isn't just convenient. It's optimal in the strong sense of being maximum likelihood. This is the deep reason OLS shows up everywhere.

\section{Properties of MLE}

Under regularity conditions (the parameter space is open, the likelihood is differentiable, the support of the data doesn't depend on the parameter), the MLE has remarkable asymptotic properties:

\paragraph{Consistency.} $\hat\theta_n \to \theta_0$ in probability as $n \to \infty$. The MLE converges to the true parameter.

\paragraph{Asymptotic normality.}
\[
  \sqrt{n}(\hat\theta_n - \theta_0) \xrightarrow{d} N\!\left(0, \mathcal{I}(\theta_0)^{-1}\right),
\]
where $\mathcal{I}(\theta) = -\E\!\left[\frac{d^2 \ell}{d\theta^2}\right]$ is the \emph{Fisher information}. The MLE is asymptotically normal with a specific (calculable) variance.

\paragraph{Asymptotic efficiency.} The MLE achieves the Cramér-Rao lower bound. No other consistent estimator can have smaller asymptotic variance. The MLE is, in this sense, optimal.

\paragraph{Invariance.} If $\hat\theta$ is the MLE of $\theta$, then $g(\hat\theta)$ is the MLE of $g(\theta)$ for any (smooth) function $g$.

These properties are why MLE is the default. It's not the only good estimator, but it has nearly all the properties you'd want.

\section{Standard errors from the Hessian}

The asymptotic variance of the MLE involves the Fisher information. In practice, you estimate it by the negative Hessian of the log-likelihood at $\hat\theta$:
\[
  \widehat{\Var}(\hat\theta) \approx \left[-\frac{d^2 \ell}{d\theta\,d\theta^\top}(\hat\theta)\right]^{-1}.
\]
This is what software packages report. R's \texttt{glm()} fits MLEs and reports standard errors from this formula. \texttt{optim()} can return the Hessian if you ask for it; you invert it to get the variance.

\section{Likelihood ratio tests}

To compare nested models, the likelihood ratio statistic is
\begin{equation}
  \text{LR} = -2\bigl(\ell_{\text{restricted}} - \ell_{\text{unrestricted}}\bigr).
  \label{eq:lr-test}
\end{equation}
Under the null hypothesis (the restriction holds), LR converges in distribution to $\chi^2_q$, where $q$ is the number of restrictions. This is one of three asymptotically equivalent tests (Wald, score, LR) for comparing models, and it's the most commonly reported.

\begin{keyinsight}
The MLE picks parameter values that maximize the probability of the observed data.

For the normal linear model, MLE is OLS. For binary outcomes, MLE is logistic regression. For count data, MLE is Poisson regression.

The MLE is consistent, asymptotically normal, and asymptotically efficient. Standard errors come from the Hessian of the log-likelihood. Likelihood ratio tests compare nested models.
\end{keyinsight}

\begin{codelab}
\begin{lstlisting}[language=Rlang]
# MLE for normal mean and variance
set.seed(2024)
data <- rnorm(100, mean = 5, sd = 3)

nll <- function(par) {
  mu <- par[1]; sigma <- par[2]
  if (sigma <= 0) return(1e10)
  -sum(dnorm(data, mu, sigma, log = TRUE))
}

fit <- optim(par = c(0, 1), fn = nll, hessian = TRUE)
fit$par                            # MLEs of mu and sigma
sqrt(diag(solve(fit$hessian)))     # standard errors

# Logistic regression as MLE
set.seed(7)
n <- 200
x <- rnorm(n)
p <- 1 / (1 + exp(-(0.5 + 1.2 * x)))
y <- rbinom(n, 1, p)

# Built-in: glm()
glm(y ~ x, family = binomial)

# Manual MLE for the same problem
log_lik <- function(beta) {
  eta <- beta[1] + beta[2] * x
  p <- 1 / (1 + exp(-eta))
  -sum(y * log(p) + (1 - y) * log(1 - p))
}
optim(c(0, 0), log_lik)$par   # should match glm()
\end{lstlisting}
\end{codelab}

\section{Practice problems}

\begin{enumerate}
  \item Show that the MLE of $p$ for $X_1, \ldots, X_n \sim \text{Bernoulli}(p)$ is $\hat p = \bar{x}$.
  \item For $X_1, \ldots, X_n \sim \text{Poisson}(\lambda)$, show that the MLE of $\lambda$ is $\bar{x}$.
  \item For an exponential distribution with rate $\lambda$ ($f(x) = \lambda e^{-\lambda x}$, $x \geq 0$), find the MLE of $\lambda$.
  \item Why is the MLE of $\sigma^2$ in the normal model $\frac{1}{n}\sum(x_i - \bar{x})^2$ rather than $\frac{1}{n-1}\sum(x_i - \bar{x})^2$? What does this say about the trade-off between MLE and unbiasedness?
  \item In R, simulate data from a Gamma distribution, then write a function for the negative log-likelihood and use \texttt{optim} to find the MLEs of the shape and rate.
  \item Prove that maximizing the likelihood is equivalent to maximizing the log-likelihood. Where does the equivalence break down?
\end{enumerate}

\begin{reflect}
$\triangleright$~State the MLE principle in your own words.

\smallskip
$\triangleright$~Why is MLE called ``maximum likelihood'' rather than ``maximum probability''?

\smallskip
$\triangleright$~Without looking back, write down the relationship between OLS and MLE for the normal linear model. Why is this the deep reason OLS is so widely used?
\end{reflect}
